\documentclass[10pt,letterpaper]{article}
\usepackage{cornell-notes}

\title{Chapter 10: Case Study 1: UNIX, Linux, and Android}
\author{}
\date{}

\setDocTitle{Chapter 10: Case Study 1: UNIX, Linux, and Android}
\setDocSubtitle{Operating Systems}
\setDocVersion{v1.0}
\setDocStatus{Study Companion}
\setDocOwner{Operating Systems Cornell Notes}
\setDocAuthor{}
\setDocDate{\today}
\setDocSummary{Cornell-format study and review notes for Chapter 10: Case Study 1: UNIX, Linux, and Android.}

\setCornellCollection{Operating Systems Cornell Notes}
\setCornellUnitType{Chapter}
\setCornellUnitNumber{10}
\setCornellUnitTitle{Case Study 1: UNIX, Linux, and Android}
\setCornellCompanionLabel{Study and review companion}

\hypersetup{pdftitle={Chapter 10: Case Study 1: UNIX, Linux, and Android},pdfsubject={Cornell notes},pdfauthor={}}

\begin{document}
\maketitle
\section*{Learning Objectives}
\begin{studybox}{By the end of this chapter, you should be able to}
\begin{enumerate}
  \item Trace the UNIX lineage through portable UNIX, Berkeley developments, standards, MINIX, and Linux.
  \item Explain the shell, library, system-call, utility, and kernel interfaces described for Linux.
  \item Trace Linux task creation, execution, scheduling, and booting.
  \item Explain Linux virtual-memory areas, page management, and relevant system calls.
  \item Trace Linux I/O, networking, driver, and module paths.
  \item Explain Linux file descriptors, open-file descriptions, i-nodes, VFS, and NFS.
  \item Analyze Linux identity, permissions, and security system calls.
  \item Explain the source-era Android architecture, Dalvik runtime, Binder IPC, application components, intents, sandboxes, security, and process model.
\end{enumerate}
\end{studybox}

\begin{studybox}{Section Map}
\hyperlink{sec-10-1}{10.1} \quad \hyperlink{sec-10-2}{10.2} \quad \hyperlink{sec-10-3}{10.3} \quad \hyperlink{sec-10-4}{10.4} \quad \hyperlink{sec-10-5}{10.5} \quad \hyperlink{sec-10-6}{10.6} \quad \hyperlink{sec-10-7}{10.7} \quad \hyperlink{sec-10-8}{10.8} \quad \hyperlink{sec-10-9}{10.9}
\end{studybox}

\section*{Cornell Cue-and-Notes Area}
\small
\begin{longtable}{C N}
\rowcolor{navy}
\color{white}\textbf{Cue / Recall Prompt} &
\color{white}\textbf{Notes}\\
\endfirsthead
\rowcolor{navy}
\color{white}\textbf{Cue / Recall Prompt} &
\color{white}\textbf{Notes (continued)}\\
\endhead
\rowcolor{ice}\multicolumn{2}{l}{\hypertarget{sec-10-1}{}\textbf{Section 10.1: History of UNIX and Linux}}\\*\hline
\textbf{10.1.1 UNICS} & \textcolor{legacy}{\textbf{Historical or legacy context:}} A small system emerged after the MULTICS project, emphasizing a simple file model, processes, a shell, and composable tools.\\\hline
\textbf{10.1.2 PDP-11 UNIX} & \textcolor{legacy}{\textbf{Historical or legacy context:}} Reimplementation on the PDP-11 and development of pipes and familiar utilities established many durable UNIX conventions.\\\hline
\textbf{10.1.3 Portable UNIX} & \textcolor{legacy}{\textbf{Historical or legacy context:}} Rewriting most of the system in C separated machine-dependent code from general kernel logic and made ports to new hardware practical.\\\hline
\textbf{10.1.4 Berkeley UNIX} & \textcolor{legacy}{\textbf{Historical or legacy context:}} BSD work added networking, virtual memory, a fast file system, shells, and utilities that influenced later UNIX systems.\\\hline
\textbf{10.1.5 Standard UNIX} & \textcolor{legacy}{\textbf{Historical or legacy context:}} POSIX and related standards specified portable library and system behavior across diverging UNIX families.\\\hline
\textbf{10.1.6 MINIX} & \textcolor{legacy}{\textbf{Historical or legacy context:}} MINIX used a small educational microkernel design with user-space servers and helped make operating-system implementation accessible for study.\\\hline
\textbf{10.1.7 Linux} & \textcolor{legacy}{\textbf{Historical or legacy context:}} Linux began as a UNIX-compatible kernel and grew through open development, GNU tools, portable interfaces, modules, and broad hardware support.\\\hline
\rowcolor{ice}\multicolumn{2}{l}{\hypertarget{sec-10-2}{}\textbf{Section 10.2: Overview of Linux}}\\*\hline
\textbf{10.2.1 Goals} & \textcolor{legacy}{\textbf{Historical or legacy context:}} The source-era goals include UNIX compatibility, efficiency, portability, multiuser protection, networking, and support from small machines through servers.\\\hline
\textbf{10.2.2 Interfaces} & Programs interact through the system-call interface, C library, shell, graphical environment, and utilities. POSIX standardizes many library-visible operations rather than a hardware-specific trap convention.\\\hline
\textbf{10.2.3 Shell} & The shell parses commands, expands names and variables, starts processes, redirects descriptors, and connects pipelines. Control structures turn the command interpreter into a programming environment.\\\hline
\textbf{10.2.4 Utilities} & Small programs perform file, text, process, compilation, administration, and communication tasks. Standard streams and pipelines allow utilities to be composed.\\\hline
\textbf{10.2.5 Kernel structure} & The Linux kernel is monolithic but modular: subsystems share a privileged address space while loadable modules and defined interfaces permit dynamic extension.\\\hline
\rowcolor{ice}\multicolumn{2}{l}{\hypertarget{sec-10-3}{}\textbf{Section 10.3: Processes in Linux}}\\*\hline
\textbf{10.3.1 Concepts} & A Linux execution entity is represented as a task. Tasks may share selected resources, allowing processes and threads to be expressed through different sharing choices.\\\hline
\textbf{10.3.2 Process calls} & \texttt{fork} creates a child using copy-on-write state; \texttt{execve} replaces the address-space image; \texttt{wait} collects child status; signals notify or control tasks; identifiers and credentials organize relationships.\\\hline
\textbf{10.3.3 Implementation} & A task structure points to scheduling, credentials, files, signal, and memory-management state. \texttt{clone} flags determine which structures are shared, and copy on write defers physical duplication.\\\hline
\textbf{10.3.4 Scheduling} & \textcolor{legacy}{\textbf{Historical or legacy context:}} The described scheduler represents runnable tasks in per-CPU structures and uses weighted fair scheduling with a red-black tree to approximate proportional CPU service.\\\hline
\textbf{10.3.5 Booting} & Firmware loads a boot loader, which loads and starts the kernel. The kernel initializes memory, interrupts, devices, and schedulers, mounts an initial root environment, and starts the first user-space process.\\\hline
\rowcolor{ice}\multicolumn{2}{l}{\hypertarget{sec-10-4}{}\textbf{Section 10.4: Memory Management in Linux}}\\*\hline
\textbf{10.4.1 Concepts} & Each process has virtual-memory areas describing contiguous regions with common backing and permissions, such as text, data, heap, stack, shared libraries, and mapped files.\\\hline
\textbf{10.4.2 System calls} & \texttt{brk} adjusts the traditional heap; \texttt{mmap} creates file or anonymous mappings; \texttt{munmap} removes them; protection and locking calls alter access or residency.\\\hline
\textbf{10.4.3 Implementation} & Memory descriptors and area structures describe address spaces; page tables map resident pages; faults allocate anonymous pages, load file data, or perform copy on write.\\\hline
\textbf{10.4.4 Paging} & \textcolor{legacy}{\textbf{Historical or legacy context:}} The described system tracks per-page state, reclaims from active and inactive sets, writes dirty pages through backing objects, and maintains free memory through background work.\\\hline
\rowcolor{ice}\multicolumn{2}{l}{\hypertarget{sec-10-5}{}\textbf{Section 10.5: Input/Output in Linux}}\\*\hline
\textbf{10.5.1 Concepts} & Devices are represented through special files, major and minor numbers, driver operations, and kernel device objects. Block and character devices follow different transfer models.\\\hline
\textbf{10.5.2 Networking} & Packets traverse protocol layers and network-device drivers rather than ordinary character or block file paths. Sockets provide the user-visible endpoint abstraction.\\\hline
\textbf{10.5.3 I/O calls} & \texttt{open}, \texttt{read}, \texttt{write}, \texttt{ioctl}, \texttt{select} or \texttt{poll}, and socket calls express data transfer, device control, readiness, and network communication.\\\hline
\textbf{10.5.4 Implementation} & The VFS and I/O layers dispatch operations through file and driver function tables, manage buffers or page cache, queue requests, and complete them through interrupts or deferred work.\\\hline
\textbf{10.5.5 Modules} & Loadable modules add drivers, file systems, or other kernel code at runtime. Symbol resolution and reference counts manage linkage and safe unloading, but modules execute with kernel privilege.\\\hline
\rowcolor{ice}\multicolumn{2}{l}{\hypertarget{sec-10-6}{}\textbf{Section 10.6: The Linux File System}}\\*\hline
\textbf{10.6.1 Concepts} & One rooted namespace contains regular files, directories, devices, links, and mounted file systems. Permissions, ownership, descriptors, and i-nodes provide access and identity.\\\hline
\textbf{10.6.2 File-system calls} & Path operations open or alter namespace objects; descriptor operations read, write, seek, duplicate, inspect, and close open files. Mounting attaches a file-system root to a directory.\\\hline
\textbf{10.6.3 Implementation} & A per-process descriptor table points to open-file descriptions, which record offsets and modes and point to VFS objects and i-nodes. The VFS dispatches to concrete file-system operations and caches data and metadata.\\\hline
\textbf{10.6.4 NFS} & NFS exposes remote directories through a request protocol and file handles. Client caching and nominally stateless servers improve performance and recovery but produce weaker consistency than a purely local file system.\\\hline
\rowcolor{ice}\multicolumn{2}{l}{\hypertarget{sec-10-7}{}\textbf{Section 10.7: Security in Linux}}\\*\hline
\textbf{10.7.1 Concepts} & Processes carry user and group identities and capabilities; files supply owner, group, mode bits, and optional richer policy. Directory execute permission controls search rather than ordinary file execution.\\\hline
\textbf{10.7.2 Security calls} & Calls change identities, groups, modes, ownership, roots, masks, and capabilities subject to current privilege. Set-user-ID execution can change effective authority.\\\hline
\textbf{10.7.3 Implementation} & Credential structures accompany tasks; access checks combine credentials, object metadata, mount rules, and security hooks. Least privilege narrows the authority inherited by services.\\\hline
\rowcolor{ice}\multicolumn{2}{l}{\hypertarget{sec-10-8}{}\textbf{Section 10.8: Android}}\\*\hline
\textbf{10.8.1 Android and Google} & \textcolor{legacy}{\textbf{Historical or legacy context:}} The described platform combines a Linux kernel with mobile middleware, application frameworks, Google-led ecosystem components, and device-vendor integration.\\\hline
\textbf{10.8.2 History} & \textcolor{legacy}{\textbf{Historical or legacy context:}} The source traces early company development, acquisition, the handset alliance, initial releases, and rapid version and device growth through the source era.\\\hline
\textbf{10.8.3 Design goals} & The platform prioritizes mobile usability, application isolation, component reuse, managed lifecycle, power awareness, device diversity, and controlled access to personal data and sensors.\\\hline
\textbf{10.8.4 Architecture} & \textcolor{legacy}{\textbf{Historical or legacy context:}} Layers include the Linux kernel, native libraries, runtime, application framework services, and applications. Framework managers expose components through defined interfaces.\\\hline
\textbf{10.8.5 Linux extensions} & \textcolor{legacy}{\textbf{Historical or legacy context:}} Mobile-specific kernel mechanisms described include Binder support, shared-memory facilities, wake locks, logging, and process or memory management adaptations.\\\hline
\textbf{10.8.6 Dalvik} & \textcolor{legacy}{\textbf{Historical or legacy context:}} Application bytecode was translated into register-oriented Dalvik form and executed in per-application runtime instances designed for constrained devices.\\\hline
\textbf{10.8.7 Binder IPC} & Binder provides protected object-oriented transactions. A kernel driver moves requests and references between processes while framework services define interfaces and permissions.\\\hline
\textbf{10.8.8 Applications} & Packages declare components and permissions in a manifest. Activities present interfaces, services perform background work, receivers handle broadcasts, and content providers share structured data.\\\hline
\textbf{10.8.9 Intents} & An intent describes an action and optional data or target. The framework resolves implicit intents against registered filters, decoupling a requester from a specific component.\\\hline
\textbf{10.8.10 Sandboxes} & Each application normally receives a distinct Linux identity and process boundary. File permissions and Binder checks isolate packages unless they deliberately share interfaces or credentials.\\\hline
\textbf{10.8.11 Security} & Permissions protect sensitive APIs and components, signatures can establish common origin, and application isolation limits direct access. User consent and platform trust remain central assumptions.\\\hline
\textbf{10.8.12 Process model} & The framework starts and stops application processes according to component demand and memory pressure. Application components must preserve state because process lifetime is not equivalent to visible component lifetime.\\\hline
\rowcolor{ice}\multicolumn{2}{l}{\hypertarget{sec-10-9}{}\textbf{Section 10.9: Summary}}\\*\hline
\textbf{Integrated case study} & Linux integrates tasks, paging, drivers, VFS, and credentials in a modular monolithic kernel. The source-era Android stack builds a managed component and IPC model on that kernel foundation.\\\hline
\end{longtable}
\normalsize

\section*{Chapter Synthesis}
\begin{studybox}{Chapter Summary}
The case study connects abstract operating-system mechanisms to concrete UNIX and Linux interfaces and data structures. Linux expresses processes and threads as tasks with selective sharing, uses paged virtual memory, dispatches I/O and file operations through function tables and VFS objects, and bases protection on credentials and object metadata. The source-era Android platform adds managed components, Binder transactions, package sandboxes, permissions, and lifecycle control.
\end{studybox}

\begin{studybox}{Key Relationships}
\begin{itemize}
  \item \texttt{fork}, copy on write, page faults, and \texttt{execve} connect process creation directly to memory management.
  \item Descriptor tables, open-file descriptions, i-nodes, and VFS operations form successive layers from process handle to persistent object.
  \item Linux credentials govern file checks and privileged operations; Android adds package identity and Binder-mediated permissions.
  \item Android component lifecycles are managed above Linux process lifetimes, so applications must tolerate process reclamation.
\end{itemize}
\end{studybox}

\begin{studybox}{Design Tradeoffs}
\begin{itemize}
  \item A modular monolithic kernel offers direct calls and broad performance but gives modules full privileged access.
  \item Copy on write makes creation cheap until writes force private copies.
  \item Client caching improves NFS performance but weakens immediate cross-client consistency.
  \item Automatic mobile process reclamation saves memory but requires explicit application state preservation.
\end{itemize}
\end{studybox}

\begin{studybox}{Common Confusions}
\begin{itemize}
  \item A Linux task may be a whole process or one thread, depending on which resources it shares.
  \item A file descriptor is per process; an open-file description may be shared and carries the current offset.
  \item A Linux file mode and an Android application permission operate at different layers and protect different interfaces.
  \item An Android component can be logically active even though the framework may later recreate its process.
\end{itemize}
\end{studybox}

\section*{Key Terms}
\small
\begin{longtable}{C N}
\rowcolor{navy}\color{white}\textbf{Term} & \color{white}\textbf{Concise definition}\\
\endfirsthead
\rowcolor{navy}\color{white}\textbf{Term} & \color{white}\textbf{Concise definition (continued)}\\
\endhead
\textbf{POSIX} & Portable interface family defining widely shared UNIX behavior.\\\hline
\textbf{Task} & Linux kernel representation of an execution entity.\\\hline
\textbf{\texttt{fork}} & Operation creating a child process from the caller's state.\\\hline
\textbf{\texttt{execve}} & Operation replacing the current process image.\\\hline
\textbf{Copy on write} & Shared read-only state copied only when modified.\\\hline
\textbf{Virtual-memory area} & Contiguous address range with common backing and permissions.\\\hline
\textbf{Loadable module} & Kernel code linked and activated at runtime.\\\hline
\textbf{VFS} & Common object and operation layer above concrete file systems.\\\hline
\textbf{Open-file description} & Kernel object containing mode, offset, and target file state.\\\hline
\textbf{i-node} & File identity and metadata object used by UNIX-style file systems.\\\hline
\textbf{NFS} & Protocol and file system for accessing remote files through mounted namespaces.\\\hline
\textbf{Effective user ID} & Credential identity used for many access checks.\\\hline
\textbf{Binder} & Android kernel-supported IPC and object-reference mechanism.\\\hline
\textbf{Intent} & Android message describing an action, data, or component request.\\\hline
\textbf{Application sandbox} & Per-package process, identity, and file isolation boundary.\\\hline
\textbf{Manifest} & Package declaration of components, permissions, and related metadata.\\\hline
\end{longtable}
\normalsize

\enlargethispage{2\baselineskip}
\begin{studybox}{Active-Recall Review}
\small
\begin{enumerate}
  \item Why did implementation in C improve UNIX portability?
  \item How do the shell, C library, and system-call interface differ?
  \item How can Linux represent both processes and threads with tasks?
  \item How do \texttt{fork}, copy on write, and \texttt{execve} interact?
  \item What sequence connects a file descriptor to a concrete file-system implementation?
  \item How do NFS caching and server state choices affect semantics?
  \item How do real and effective identities influence access?
  \item What role does Binder play between Android application processes?
  \item How do activities, services, receivers, and content providers differ?
  \item Why must an Android application separate component state from process lifetime?
\end{enumerate}
\end{studybox}

\par\medskip
{\color{navy}\bfseries Next-Chapter Connections}\par\smallskip
Chapter 11 provides a contrasting case study organized around Windows objects, handles, executive subsystems, priority scheduling, working sets, I/O request packets, NTFS, and access tokens.
\normalsize
\end{document}
